\documentclass{article}
\usepackage[margin=2cm]{geometry}
\usepackage{amsmath}
\usepackage{graphicx}
\usepackage{booktabs}
\usepackage[utf8]{inputenc}
\begin{document}
\title{
1 Gamma matrix and spinor conventions
}
For concreteness, we take the following basis of gamma matrices
\[
\begin{array}{l}
\gamma_{1}=\sigma_{2} \otimes \mathbf{1}_{2} \otimes \sigma_{3} \\
\gamma_{2}=\sigma_{2} \otimes \mathbf{1}_{2} \ot� \sigma_{1} \\
\gamma_{3}=\mathbf{1}_{2} \otimes \sigma_{1} \otimes \sigma_{2} \\
\gamma_{4}=\mathbf{1}_{2} \otimes \sigma_{3} \ot� \sigma_{2} \\
\gamma_{5}=\sigma_{1} \otimes \sigma_{2} \otimes \mathbf{1}_{2} \\
\gamma_{6}=\sigma_{3} \otimes \sigma_{2} \otimes \mathbf{1}_{2}
\end{array}
\]
These gamma matrices satisfy the Clifford algebra
\[
\left\{\gamma_{\mu}, \gamma_{\nu}\right\}=2 \delta_{\mu \nu}
\]
as appropriate for a positive definite Euclidean spacetime. All matrices are purely imaginary and satisfy
\[
\left(\gamma_{\mu}\right)^{\dagger}=\gamma_{\mu} \quad\left(\gamma_{\mu}\right)^{2}=\mathbf{1}
\]
We will now be interested in a seven-dimensional Clifford algebra, which will require the introduction of a new matrix \(\gamma_{7}\). The reason we are interested in this is that we would like to represent hyperbolic space \(\mathbb{H}_{6}\) as a hypersurface in a seven-dimensional ambient space. This allows us to determine properties of the Dirac spinors in the Euclidean-continued \(F(4)\) gauged supergravity theory with \(\mathbb{H}_{6}\) background by first considering Dirac spinors in seven dimensions and then performing a timelike reduction. In particular, we will choose a 7D metric of signature \((+,++,+,+,+,+,-)\) for the ambient space. Then hyperbolic space \(\mathbb{H}_{6}\) is given by the following quadratic form
\[
x_{1}^{2}+\cdots+x_{6}^{2}-x_{7}^{2}=-L^{2}
\]
The seven-dimensional Clifford algebra is made up of the set of matrices \(\left\{\gamma_{1}, \ldots, \gamma_{6}, \gamma_{7}\right\}\), with \(\gamma_{7}\) satisfying
\[
\left(\gamma_{7}\right)^{2}=-\mathbf{1} \quad\left\{\gamma_{\mu}, \gamma_{7}\right\}=0 \quad \forall \mu \neq 7
\]
As usual, we use the notation \(\gamma^{7}=(\gamma_{7})^{-1}\), so that by the above we have \(\gamma^{7}=-\gamma_{7}\). We now discuss Dirac spinors in \(d=7\). We define the Dirac conjugate of \(\psi_{A}\) to be
\[
\bar{\psi}_{A}=\psi_{A}^{\dagger} G^{-1}
\]
for some matrix \(G\). There are two possible choices for \(G\), which in the particular case of the ambient space above are
\[
G_{1}=\gamma^{7} \quad G_{2}=\gamma^{1} \cdots \gamma^{6}
\]
\end{document}
